\documentclass[11 pt]{article}
\usepackage[a4paper, margin = 1 in]{geometry}
\usepackage[spanish]{babel}
\usepackage{minted}
\usepackage{amsmath}
\usepackage{amssymb}

\title{Programación Competitiva}
\author{Arturo Caballero Ortega}

\begin{document}

\begin{titlepage}

    \centering
    \vspace*{\fill}
    {\Huge \textbf{Programación Competitiva} \par}
    \vspace*{\fill}

    {\Large Por: Arturo Caballero Ortega\par}
    \vspace{1cm}
    {\Large \today\par}
    \vspace{2cm}

\end{titlepage}

\tableofcontents
\newpage

\section{Búsqueda binaria}


A continuación un ejemplo de búsqueda binaria sobre arrays:


\begin{minted}[frame=single, fontsize=\small, linenos]{cpp}
int a[ll] = {/* array ordenado */};

int l = 0, r = N;
while(l < r){
    int m = (l + r) / 2;
    if(predicado(a[m])) l = m;
    else r = m;
}

// l es el índice del primer elemento que cumple predicado
\end{minted}

Para el caso de trabajar con números reales, nos puede bastar con comprobar, en vez de la condición
$l<r$, la condición $l + \varepsilon < r$, donde $\varepsilon$ es un valor muy pequeño, $1e9$ por ejemplo.

\vspace{25 px}
\section{Optimizaciones básicas}

    
Destaquemos:
\begin{itemize}
    \item Two pointers
    \item Sliding window
    \item Counting sort
    \item Majority element
\end{itemize}

TODO: Explicarlos brevemente.

\vspace{25 px}
\section{Tecnicas de ordenación}


Destacaremos tres problemas de ordenación.

\subsection{Problema 1}
En el primer problema tenemos una permutación de $1,2,\dots ,N$ y podemos hacer swaps de elementos consecutivos.
El objetivo es ordenar la permutación en el menor número de swaps posible, obtener al menos el número de swaps necesarios.


En este caso, el número de swaps necesarios es el mismo número de inversiones en la permutación. Una
inversión es un par de elementos $a_i$ y $a_j$ tales que $i < j$ y $a_i > a_j$.\\
Para obtener el número de inversiones podemos utilizar un algoritmo basado en mergesort, tal que:

\begin{minted}[frame=single, fontsize=\small, linenos]{cpp}
vector<int> mergesort(vector<int> & v, int l, int r){
    if(l == r){
        return {v[l]};
    }
    int m = (l+r)/2;
    vector<int> a = mergesort(v, l, m);
    vector<int> b = mergesort(v, m+1, r);
    vector<int> nuevo;
    auto it1 = a.begin(), it2 = b.begin();
    while(it1 != a.end() && it2 != b.end()){
        if(*it1 < *it2){
            nuevo.push_back(*it1);
            it1++;
            continue;
        }
        nuevo.push_back(*it2);
        inversiones += (a.end() - it1);
        it2++;
    }
    while(it1 != a.end()){
        nuevo.push_back(*it1);
        it1++;
    }
    while(it2 != b.end()){
        nuevo.push_back(*it2);
        it2++;
    }
    return nuevo;
}
\end{minted}

\subsection{Problema 2}
En el segundo problema tenemos que obtener la mayor subsecuencia creciente de un array $v$ de tamaño $N$.
Para este problema necesitamos definir un array auxiliar $a$. Recorreremos el array $v$ con un algoritmo
similar al siguiente:

\begin{minted}[frame=single, fontsize=\small, linenos]{cpp}
for(int i=0; i<n; ++i){
    auto it = lower_bound(a.begin(), a.end(), v[i]);
    if(it == a.end()) {
        a.push_back(v[i]);
    }
    else {
        *it = v[i];
    }
}
// La longitud de la mayor subsecuencia creciente es a.size()
\end{minted}

\subsection{Problema 3}
En el tercer problema tenemos que ordenar una permutación con la posibilidad de hacer swaps de dos elementos cualesquiera. El objetivo es obtener el menor número de swaps posibles.\\
En este problema tenemos que ver la permutación como un grafo, donde el elemento en la posición $i$ apunta al elemento en la posición $a_i$. Con esta representación, cada grafo conexo formado de tamaño $k$ necesita $k-1$ swaps para ser ordenado.


De hecho, sea $n$ el número de elementos, $j$ el número de ciclos y $k_i$ el tamaño del ciclo $i$-ésimo.
Entonces el número de swaps necesarios es $n-j$:
\[ \sum_{i=1}^{j} \left( k_i - 1 \right) = \sum_{i=1}^{j} k_i - j = n - j \]

\vspace{25 px}
\section{Caminos mínimos}
\subsection{Algoritmo de Dijkstra}

\begin{minted}[frame=single, fontsize=\small, linenos]{cpp}
for (int i = 1; i <= n; i++) distance[i] = INF;
distance[x] = 0;
q.push({0,x});
while (!q.empty()) {
    int a = q.top().second; q.pop();
    if (processed[a]) continue;
    processed[a] = true;
    for (auto u : adj[a]) {
        int b = u.first, w = u.second;
        if (distance[a]+w < distance[b]) {
            distance[b] = distance[a]+w;
            q.push({-distance[b],b});
        }
    }
}
\end{minted}


\subsection{Algoritmo de Floyd-Warshall}

\begin{minted}[frame=single, fontsize=\small, linenos]{cpp}
for (int i = 1; i <= n; i++) {
    for (int j = 1; j <= n; j++) {
        if (i == j) distance[i][j] = 0;
        else if (adj[i][j]) distance[i][j] = adj[i][j];
        else distance[i][j] = INF;
    }
}
\end{minted}


\begin{minted}[frame=single, fontsize=\small, linenos]{cpp}
for (int k = 1; k <= n; k++) {
    for (int i = 1; i <= n; i++) {
        for (int j = 1; j <= n; j++) {
            distance[i][j] = min(distance[i][j],
            distance[i][k]+distance[k][j]);
        }
    }
}
\end{minted}

\vspace{25 px}
\section{Matemáticas}
\subsection{Criba de Eratóstenes}

\begin{minted}[frame=single, fontsize=\small, linenos]{cpp}
int n;
vector<bool> is_prime(n + 1, true);
is_prime[0] = is_prime[1] = false;
for (int i = 2; i * i <= n; i++) {
    if (is_prime[i]) {
        for (int j = i * i; j <= n; j += i)
            is_prime[j] = false;
    }
}
\end{minted}

Para obtener la factorización de un número se puede hacer algo similar a esto:
\begin{minted}[frame=single, fontsize=\small, linenos]{cpp}
vector<int> max_prime (n+1);
for(int i=0; i<n+1; ++i) max_prime[i] = i;
for (int i = 2; i * i <= n; i++) {
    if (max_prime [i] == i) {
        for (int j = i + i; j <= n; j += i)
        max_prime[j] = i;
    }
}

vector<list<int>> a(n+1);
for(int i=2; i<n+1; ++i){
    vector <pair <int, int>> fact;
    int n = i;
    while(n != 1){
        int p = max_prime[n];
        n /= p;
        if(fact.empty () || fact.back ().first != p)
            fact.emplace_back (p, 1);
        else fact.back ().second ++;
    }

    for(auto p : fact) a[i].push_back(p.first);
}
\end{minted}

\subsection{Aritmética modular}
Algunas cosas interesantes:

\[ (x+y) \% m = ((x \% m) + (y \% m)) \% m \]
\[ (x-y) \% m = ((x \% m) - (y \% m) + m) \% m \]
\[ (x\cdot y) \% m = ((x \% m) \cdot (y \% m)) \% m \]
\[ x^n \% m = ((x \% m)^n) \% m \]

También podemos obtener el inverso modular:
\[ \left( \frac{x}{y} \right) \% m = \left( x\%m \cdot y^{m-2} \% m \right) \text{con m primo}\]

\subsection{Exponenciación rápida}

\begin{minted}[frame=single, fontsize=\small, linenos]{cpp}
int (int a, int e, int m){
    int r = 1;
    while(e){
        if(e & 1) r = (r*a) % m;
        e >>= 1;
        a = (a*a) % m;
    }
    return r;
}
\end{minted}

\subsection{Identidad de Bezout}
Supongamos no necesitar esto gracias.

\begin{minted}[frame=single, fontsize=\small, linenos]{cpp}
constexpr array<ll, 3> egcd(ll a, ll b) {
    ll ga = 1, g = a;
    ll za = 0, z = b;
    ll t, aux;
    while (z != 0) {
        aux = z;
        t = g / z;
        z = g % z;
        g = aux;
        aux = za;
        za = ga - t * za;
        ga = aux;
    }
    ll gb = b == 0 ? 0 : (g - ga * a) / b;
    return {ga, gb, g};
}
\end{minted}

\subsection{Combinatoria}

Número de subconjuntos de $k$ elementos tomados de un conjunto de $n$ elementos.
\[ \binom{n}{k} = \frac{n!}{k!(n-k)!}, \hspace{15 px} 
\binom{n}{0} = \binom{n}{n} = 1, \hspace{15 px} \binom{n}{k} = 
\binom{n-1}{k-1} + \binom{n-1}{k} \] 

\begin{minted}[frame=single, fontsize=\small, linenos]{cpp}
const int maxn = 100; // Replace 100 with the desired maximum value
int C[maxn + 1][maxn + 1];
C[0][0] = 1;
for (int n = 1; n <= maxn; ++n) {
    C[n][0] = C[n][n] = 1;
    for (int k = 1; k < n; ++k)
        C[n][k] = C[n - 1][k - 1] + C[n - 1][k];
}
\end{minted}

\subsection{Principio de inclusión-exclusión}

\[ \left\lvert \bigcup_{i=1}^n A_i \right\rvert = \sum_{i=1}^{n} \left\lvert A_i \right\rvert - \sum_{0 \le i < j \le n} \left\lvert A_i \cap A_j \right\rvert + \sum_{0 \le i < j < k \le n} \left\lvert A_i \cap A_j \cap A_k \right\rvert - \cdots + (-1)^{n-1} \left\lvert A_1 \cap A_2 \cap \cdots \cap A_n \right\rvert
\]

\end{document}